---
title: "Random Walks on GL₂(𝔽ₚ)"
author: "Eugen Nesbakken"
date: "2026-02-25"
description: "A detailed treatment of random walks on the general linear group over a finite field: group structure, generating sets, total variation distance, spectral gap, and the non-abelian Fourier-analytic upper bound."
category: mathematics
tags: [algebra, probability, representation-theory, random-walks]
---

\section{Introduction}

Let $G$ be a finite group and $\mu$ a probability measure on $G$.
The \textbf{random walk} driven by $\mu$ is the discrete-time Markov chain
\[
  X_0 = e, \qquad X_{t+1} = X_t \cdot S_t, \quad S_t \overset{\mathrm{iid}}{\sim} \mu,
\]
where $e$ is the identity element.
After $t$ steps the position $X_t$ has distribution $\mu^{*t}$, the $t$-fold
convolution of $\mu$ with itself:
\[
  \mu^{*t}(g) = \sum_{h \in G} \mu^{*(t-1)}(h)\,\mu(h^{-1}g).
\]
The fundamental question is: how large must $t$ be so that $\mu^{*t}$ is
indistinguishable from the uniform measure $U_G(g) = 1/|G|$?

We make this precise with the \textbf{total variation distance}
\[
  \|\mu^{*t} - U_G\|_{\mathrm{TV}}
  = \frac{1}{2}\sum_{g \in G}\bigl|\mu^{*t}(g) - U_G(g)\bigr|
  = \max_{A \subseteq G}\bigl(\mu^{*t}(A) - U_G(A)\bigr).
\]
We study this problem for $G = \GL_2(\mathbb{F}_p)$, the group of $2 \times 2$
invertible matrices over the finite field $\mathbb{F}_p$.

\section{The Group \texorpdfstring{$\GL_2(\mathbb{F}_p)$}{GL2(Fp)}}

\subsection{Definition and Basic Properties}

\begin{definition}[$\GL_2(\mathbb{F}_p)$]
For a prime $p$, the \textbf{general linear group} $\GL_2(\mathbb{F}_p)$ is the
set of $2 \times 2$ matrices with entries in $\mathbb{F}_p = \mathbb{Z}/p\mathbb{Z}$
having non-zero determinant, equipped with matrix multiplication modulo $p$.
\end{definition}

We represent an element as a flat vector $[a, b, c, d]$ corresponding to
$\bigl[\begin{smallmatrix}a & b \\ c & d\end{smallmatrix}\bigr]$.
The group operation is matrix multiplication modulo $p$, and the inverse is
given by the classical $2 \times 2$ adjugate formula:
\[
  \begin{pmatrix} a & b \\ c & d \end{pmatrix}^{-1}
  = (\det A)^{-1} \begin{pmatrix} d & -b \\ -c & a \end{pmatrix} \pmod{p}.
\]

\begin{lemma}[Order of $\GL_2(\mathbb{F}_p)$]
\[
  |\GL_2(\mathbb{F}_p)| = (p^2 - 1)(p^2 - p).
\]
\end{lemma}

\begin{proof}
An element of $\GL_2(\mathbb{F}_p)$ is a pair of linearly independent column
vectors in $\mathbb{F}_p^2$.
The first column may be any non-zero vector: $p^2 - 1$ choices.
Given the first column $v_1$, the second column must be linearly independent of
$v_1$, so it can be any vector except the $p$ scalar multiples $\lambda v_1$
($\lambda \in \mathbb{F}_p$): $p^2 - p$ choices.
All pairs satisfying this condition yield an invertible matrix, giving
$(p^2 - 1)(p^2 - p)$ elements in total.
\end{proof}

This can be rewritten as $p(p-1)^2(p+1)$, which grows like $p^4$.
The first few values are:
\[
  |\GL_2(\mathbb{F}_2)| = 6, \quad
  |\GL_2(\mathbb{F}_3)| = 48, \quad
  |\GL_2(\mathbb{F}_5)| = 480, \quad
  |\GL_2(\mathbb{F}_7)| = 2016.
\]

\subsection{The Subgroup \texorpdfstring{$\SL_2(\mathbb{F}_p)$}{SL2(Fp)}}

The determinant map $\det : \GL_2(\mathbb{F}_p) \to \mathbb{F}_p^\times$ is a
surjective group homomorphism.
Its kernel is the \textbf{special linear group}
\[
  \SL_2(\mathbb{F}_p) = \ker(\det) = \{A \in \GL_2(\mathbb{F}_p) : \det A = 1\}.
\]
By the first isomorphism theorem,
$\GL_2(\mathbb{F}_p) / \SL_2(\mathbb{F}_p) \cong \mathbb{F}_p^\times$,
so $\SL_2$ is a normal subgroup of index $p - 1$ and order
\[
  |\SL_2(\mathbb{F}_p)| = \frac{(p^2-1)(p^2-p)}{p-1} = p(p^2-1).
\]
This subgroup plays a crucial role in understanding the generating sets for our
random walks.

\subsection{Conjugacy Classes}

Every element of $\GL_2(\mathbb{F}_p)$ is conjugate to exactly one of the
following canonical forms, determined by its rational canonical (Jordan) form
over $\mathbb{F}_p$:

\begin{center}
\renewcommand{\arraystretch}{1.4}
\begin{tabular}{llll}
\hline
\textbf{Type} & \textbf{Representative} & \textbf{\# classes} & \textbf{Class size} \\
\hline
Scalar & $\lambda I$ & $p - 1$ & $1$ \\
Jordan (non-scalar) & $\bigl[\begin{smallmatrix}\lambda & 1 \\ 0 & \lambda\end{smallmatrix}\bigr]$ & $p - 1$ & $p^2 - 1$ \\
Split semisimple & $\bigl[\begin{smallmatrix}\lambda & 0 \\ 0 & \mu\end{smallmatrix}\bigr]$, $\lambda \neq \mu$ & $\binom{p-1}{2}$ & $p(p+1)$ \\
Non-split semisimple & companion of irred.\ $x^2 - tx + n$ & $\frac{p(p-1)}{2}$ & $p(p-1)$ \\
\hline
\end{tabular}
\end{center}

The total number of conjugacy classes is $p^2 - 1$, consistent with the
general fact that the number of irreducible representations equals the number of
conjugacy classes.
One verifies the class-size formula by summing:
$(p-1)\cdot 1 + (p-1)(p^2-1) + \binom{p-1}{2}p(p+1) + \frac{p(p-1)}{2}p(p-1) = (p^2-1)(p^2-p)$. $\checkmark$

The \textbf{non-split semisimple} classes correspond to elements whose
characteristic polynomial $x^2 - tx + n$ is \emph{irreducible} over
$\mathbb{F}_p$, i.e.\ the discriminant $t^2 - 4n$ is a quadratic non-residue
modulo $p$.
Their eigenvalues live in $\mathbb{F}_{p^2} \setminus \mathbb{F}_p$.

\section{Generating Sets}

\subsection{Transvections}

\begin{definition}[Transvection]
A \textbf{transvection} is an elementary matrix of the form $I + \lambda E_{ij}$
for $i \neq j$ and $\lambda \neq 0$, where $E_{ij}$ is the matrix unit with a
$1$ in position $(i,j)$ and $0$s elsewhere.
\end{definition}

For $\GL_2(\mathbb{F}_p)$ the only transvections are scalar multiples of
$T_{12} = I + E_{12}$ and $T_{21} = I + E_{21}$, i.e.
\[
  T_{12}^{(\lambda)} = \begin{pmatrix} 1 & \lambda \\ 0 & 1 \end{pmatrix}, \quad
  T_{21}^{(\lambda)} = \begin{pmatrix} 1 & 0 \\ \lambda & 1 \end{pmatrix},
  \qquad \lambda \in \mathbb{F}_p^\times.
\]
Every transvection has determinant 1 (it is unipotent), so transvections lie in
$\SL_2(\mathbb{F}_p)$.
It is a classical theorem that transvections generate $\SL_2(\mathbb{F}_p)$:

\begin{theorem}
$\SL_2(\mathbb{F}_p)$ is generated by $\{T_{12}^{(\lambda)}, T_{21}^{(\lambda)} : \lambda \in \mathbb{F}_p^\times\}$.
\end{theorem}

In our random walk we use the symmetric generating set
$S_{\mathrm{tv}} = \{T_{12}, T_{12}^{-1}, T_{21}, T_{21}^{-1}\}$
where $T_{ij} = T_{ij}^{(1)}$.
Note that $T_{12}^{-1} = T_{12}^{(p-1)}$ and $T_{21}^{-1} = T_{21}^{(p-1)}$.

\subsection{The Swap Generator and Ergodicity on $\GL_2$}

Since all transvections lie in $\SL_2$, the walk driven by $S_{\mathrm{tv}}$
is confined to $\SL_2(\mathbb{F}_p)$ — the full group $\GL_2(\mathbb{F}_p)$ is
never reached.
To make a walk that is ergodic on all of $\GL_2(\mathbb{F}_p)$, we need at
least one generator of determinant $\neq 1$.

The \textbf{swap matrix}
\[
  P = \begin{pmatrix} 0 & 1 \\ 1 & 0 \end{pmatrix}, \qquad \det P = -1,
\]
is self-inverse ($P^2 = I$) and has determinant $-1 \equiv p - 1 \pmod{p}$.
Adding $P$ to the generating set gives
$S_{\mathrm{ts}} = S_{\mathrm{tv}} \cup \{P\}$.
Since $\det P = -1$ generates $\mathbb{F}_p^\times / (\mathbb{F}_p^\times)^2$
(for $p$ odd) or equals $1$ (for $p = 2$, but $p = 2$ is special anyway),
the walk driven by $S_{\mathrm{ts}}$ reaches all elements of $\GL_2(\mathbb{F}_p)$.

\subsection{The Special Case $p = 2$: Periodicity}

Over $\mathbb{F}_2$ we have $-1 \equiv 1$, so $T_{12}^{-1} = T_{12}$ and
$T_{21}^{-1} = T_{21}$.
Both generators have order 2: $T_{12}^2 = I$ and $T_{21}^2 = I$.
Consequently the Cayley graph $\mathrm{Cay}(\GL_2(\mathbb{F}_2), S_{\mathrm{tv}})$
is \textbf{bipartite}: the vertices split into two classes $V_0, V_1$ such that
every edge goes from $V_0$ to $V_1$.
This means the walk alternates between $V_0$ and $V_1$ at every step,
and $\mu^{*t}$ never converges to uniform — the total variation distance
oscillates and is bounded below by $1/2$.

The standard remedy is a \textbf{lazy walk}: replace $\mu$ by the measure
\[
  \tilde\mu = \frac{1}{2}\delta_e + \frac{1}{2}\mu,
\]
which with probability $1/2$ holds in place and with probability $1/2$ takes a
random step.
Equivalently, one adds $e$ to the generating multiset.
The lazy walk is aperiodic (since $\tilde\mu(e) > 0$), so $\tilde\mu^{*t}$
converges to uniform.

Note that $\GL_2(\mathbb{F}_2) \cong S_3$, the symmetric group on 3 letters —
the smallest non-abelian group.

\section{Random Walks: Formal Setup}

\subsection{The Markov Chain}

Fix a symmetric probability measure $\mu$ on $G$ (i.e.\ $\mu(g) = \mu(g^{-1})$
for all $g$).
The random walk is the time-homogeneous Markov chain on $G$ with transition
probabilities $P(g, h) = \mu(g^{-1}h)$.
The \textbf{stationary distribution} is the uniform measure $U_G$, since
\[
  \sum_{g \in G} U_G(g)\,P(g, h)
  = \frac{1}{|G|}\sum_{g \in G}\mu(g^{-1}h)
  = \frac{1}{|G|}\sum_{x \in G}\mu(x)
  = \frac{1}{|G|}
  = U_G(h).
\]
Ergodicity (irreducibility + aperiodicity) is equivalent to the generating set
$\mathrm{supp}(\mu)$ generating $G$ and the lazy condition $\mu(e) > 0$
(or the Cayley graph being non-bipartite).

\subsection{The Convolution Algebra}

The space $L^2(G) = \{f : G \to \mathbb{C}\}$ is a $|G|$-dimensional Hilbert
space with inner product $\langle f, h \rangle = \sum_{g \in G} f(g)\overline{h(g)}$.
Define convolution by
\[
  (f * h)(g) = \sum_{x \in G} f(x)\,h(x^{-1}g).
\]
Then $\mu^{*t} = \mu * \mu^{*(t-1)}$, and the random walk operator
\[
  T : L^2(G) \to L^2(G), \qquad (Tf)(g) = \sum_{h \in G} \mu(h)\,f(gh) = (\mu * f)(g)
\]
is self-adjoint (since $\mu$ is symmetric) with largest eigenvalue $1$
(eigenvector: constant functions).

\section{Total Variation Distance and Convergence}

\subsection{Basic Properties of TVD}

The total variation distance satisfies:
\begin{itemize}
  \item $0 \leq \|\mu^{*t} - U_G\|_{\mathrm{TV}} \leq 1$.
  \item $\|\mu^{*(t+1)} - U_G\|_{\mathrm{TV}} \leq \|\mu^{*t} - U_G\|_{\mathrm{TV}}$
        (non-increasing in $t$, by the data-processing inequality).
  \item $\|\mu^{*t} - U_G\|_{\mathrm{TV}} \to 0$ as $t \to \infty$
        if and only if the walk is ergodic.
\end{itemize}

\subsection{The \texorpdfstring{$\varepsilon$}{epsilon}-Mixing Time}

\begin{definition}[Mixing time]
For $\varepsilon > 0$, the \textbf{$\varepsilon$-mixing time} is
\[
  t_{\mathrm{mix}}(\varepsilon) = \min\{t \geq 0 : \|\mu^{*t} - U_G\|_{\mathrm{TV}} \leq \varepsilon\}.
\]
The \textbf{mixing time} is $t_{\mathrm{mix}} = t_{\mathrm{mix}}(1/4)$ (the choice $1/4$ is conventional).
\end{definition}

\section{Spectral Theory of the Random Walk Operator}

\subsection{Eigenvalues and the Spectral Gap}

Since $T$ is a self-adjoint operator on $L^2(G)$, it has real eigenvalues.
Order them as
\[
  1 = \lambda_1 \geq \lambda_2 \geq \cdots \geq \lambda_{|G|} \geq -1.
\]
The largest eigenvalue is $\lambda_1 = 1$ (constant eigenfunctions).
For an ergodic walk, $\lambda_2 < 1$ and $\lambda_{|G|} > -1$.

\begin{definition}[Spectral gap]
The \textbf{spectral gap} is $\gamma = 1 - \lambda_2$.
The \textbf{absolute spectral gap} is $\gamma^* = 1 - \max(|\lambda_2|, |\lambda_{|G|}|)$.
\end{definition}

\subsection{The Spectral Bound on TVD}

\begin{theorem}[Spectral upper bound]
\[
  \|\mu^{*t} - U_G\|_{\mathrm{TV}} \leq \frac{\sqrt{|G|}}{2}\,(1 - \gamma^*)^t.
\]
\end{theorem}

\begin{proof}
Write $\mu^{*t} - U_G$ in the orthonormal eigenbasis $\{f_1, \ldots, f_{|G|}\}$ of $T$,
where $f_1 \equiv |G|^{-1/2}$ (the constant eigenfunction for $\lambda_1 = 1$).
Since $U_G$ is the projection onto $f_1$,
\[
  \mu^{*t}(g) - U_G(g) = \sum_{k=2}^{|G|} \lambda_k^t \langle \mu^{*t}, f_k \rangle f_k(g).
\]
By Cauchy-Schwarz and Parseval,
\begin{align*}
  2\|\mu^{*t} - U_G\|_{\mathrm{TV}}
  &= \sum_g |\mu^{*t}(g) - U_G(g)|
  \leq \sqrt{|G|} \left(\sum_g |\mu^{*t}(g) - U_G(g)|^2\right)^{1/2} \\
  &= \sqrt{|G|} \left(\sum_{k=2}^{|G|} \lambda_k^{2t}\right)^{1/2}
  \leq \sqrt{|G|}\,(1-\gamma^*)^t \left(\sum_{k=2}^{|G|} 1\right)^{1/2}
  \leq |G|\,(1-\gamma^*)^t,
\end{align*}
and a sharper version using $\|\mu^{*(t)} - U_G\|_2^2 \leq (1-\gamma^*)^{2t}$
yields the stated bound.
\end{proof}

This gives the mixing time estimate
\[
  t_{\mathrm{mix}}(\varepsilon) \leq \left\lceil \frac{\log(|G|/4\varepsilon^2)}{2\gamma^*} \right\rceil.
\]
Hence \emph{a large spectral gap means fast mixing}.

\subsection{Eigenvalues for $\GL_2(\mathbb{F}_2)$}

For $G = \GL_2(\mathbb{F}_2) \cong S_3$ with the lazy transvection generating
set $\tilde{S} = \{e, T_{12}, T_{21}\}$ (weights $1/3$ each), one can compute
the $6 \times 6$ transition matrix explicitly.
The eigenvalues are (sorted descending):
\[
  1,\; \tfrac{1}{3},\; \tfrac{1}{3},\; \tfrac{1}{3},\; -\tfrac{1}{3},\; -\tfrac{1}{3}.
\]
The spectral gap is $\gamma = 1 - 1/3 = 2/3$, giving the bound
\[
  \|\tilde\mu^{*t} - U_G\|_{\mathrm{TV}} \leq \frac{\sqrt{6}}{2}\left(\frac{1}{3}\right)^t.
\]
This decays geometrically, and the exact TVD (computed by convolution) confirms
it reaches $< 0.01$ by $t = 10$.

\section{Non-Abelian Fourier Analysis}

For a non-abelian group, the spectral bound can be refined using the full
representation theory.

\subsection{Irreducible Representations and Characters}

Recall that for a finite group $G$, the irreducible complex representations
$\rho : G \to \GL(V_\rho)$ are finite in number, and their characters
$\chi_\rho(g) = \mathrm{tr}(\rho(g))$ satisfy the orthogonality relations
\[
  \frac{1}{|G|} \sum_{g \in G} \chi_\rho(g)\,\overline{\chi_\sigma(g)} = \delta_{\rho,\sigma}.
\]
The number of irreducible representations equals the number of conjugacy
classes, so $\GL_2(\mathbb{F}_p)$ has $p^2 - 1$ irreducible representations.

\subsection{The Matrix Fourier Transform}

\begin{definition}[Matrix Fourier transform]
For $f : G \to \mathbb{C}$ and an irreducible representation $\rho$ of $G$,
define the \textbf{Fourier transform} of $f$ at $\rho$ by
\[
  \hat{f}(\rho) = \sum_{g \in G} f(g)\,\rho(g) \in \mathrm{End}(V_\rho).
\]
\end{definition}

The Fourier transform satisfies $\widehat{f * h}(\rho) = \hat{f}(\rho)\,\hat{h}(\rho)$
(convolution becomes matrix multiplication).
The \textbf{Plancherel formula} reads
\[
  \sum_{g \in G} |f(g)|^2 = \frac{1}{|G|} \sum_{\rho \in \hat{G}} d_\rho\,\|\hat{f}(\rho)\|_F^2,
\]
where $d_\rho = \dim V_\rho$ and $\|\cdot\|_F$ is the Frobenius norm.

\subsection{The Upper Bound Lemma}

\begin{theorem}[Diaconis--Shahshahani upper bound lemma, 1981]
Let $\mu$ be a symmetric probability measure on a finite group $G$.
Then
\[
  4\|\mu^{*t} - U_G\|_{\mathrm{TV}}^2
  \leq \sum_{\rho \neq \mathbf{1}} d_\rho\,\|\hat\mu(\rho)\|_F^{2t},
\]
where the sum is over all non-trivial irreducible representations $\rho$.
\end{theorem}

\begin{proof}[Proof sketch]
By Cauchy-Schwarz and the Plancherel formula,
\begin{align*}
  4\|\mu^{*t} - U_G\|_{\mathrm{TV}}^2
  &\leq |G| \sum_{g \in G}|\mu^{*t}(g) - U_G(g)|^2 \\
  &= \sum_{\rho \neq \mathbf{1}} d_\rho \|\widehat{\mu^{*t}}(\rho)\|_F^2
  = \sum_{\rho \neq \mathbf{1}} d_\rho \|\hat\mu(\rho)^t\|_F^2
  \leq \sum_{\rho \neq \mathbf{1}} d_\rho \|\hat\mu(\rho)\|_F^{2t}.
\end{align*}
The last inequality uses submultiplicativity of the Frobenius norm and the
fact that $\|\hat\mu(\rho)\|_{\mathrm{op}} \leq 1$.
\end{proof}

To apply this, one computes $\hat\mu(\rho)$ for each non-trivial irreducible $\rho$.
For the uniform measure on a generating set $S$,
$\hat\mu(\rho) = \frac{1}{|S|}\sum_{s \in S} \rho(s)$,
and the bound becomes useful once the operator norms $\|\hat\mu(\rho)\|_{\mathrm{op}} < 1$
are understood.

\section{Irreducible Representations of \texorpdfstring{$\GL_2(\mathbb{F}_p)$}{GL2(Fp)}}

The $p^2 - 1$ irreducible representations of $\GL_2(\mathbb{F}_p)$ fall into
four families, reflecting the four conjugacy class types.

\subsection{Characters of $\mathbb{F}_p^\times$ and the Determinant}

Let $\omega : \mathbb{F}_p^\times \to \mathbb{C}^\times$ be a character (group
homomorphism).
Since $\mathbb{F}_p^\times \cong \mathbb{Z}/(p-1)\mathbb{Z}$, there are $p - 1$
such characters.
Composing with $\det$ gives a one-dimensional representation $\omega \circ \det$
of $\GL_2(\mathbb{F}_p)$.
These account for $p - 1$ of the irreducibles.

\subsection{Principal Series Representations}

For characters $\chi_1, \chi_2 : \mathbb{F}_p^\times \to \mathbb{C}^\times$ with
$\chi_1 \neq \chi_2$, the \textbf{principal series representation}
$\pi(\chi_1, \chi_2)$ is obtained by parabolic induction from the Borel
subgroup of upper-triangular matrices.
It has dimension $p + 1$.
Since $\pi(\chi_1, \chi_2) \cong \pi(\chi_2, \chi_1)$, there are
$\binom{p-1}{2}$ such representations.

\subsection{Special (Steinberg) Representations}

For each character $\chi$ of $\mathbb{F}_p^\times$, there is a
\textbf{special representation} $\mathrm{St} \otimes (\chi \circ \det)$ of
dimension $p$, where $\mathrm{St}$ is the Steinberg representation.
This gives $p - 1$ representations of dimension $p$.

\subsection{Cuspidal Representations}

The remaining $\frac{p(p-1)}{2}$ irreducibles are \textbf{cuspidal representations},
each of dimension $p - 1$.
They are constructed via induction from characters of $\mathbb{F}_{p^2}^\times$
that do not factor through the norm map $N : \mathbb{F}_{p^2}^\times \to \mathbb{F}_p^\times$.
Cuspidal representations are the most subtle; they do not appear in any
principal series.

\subsection{Dimension Check}

One verifies that
\[
  (p-1) \cdot 1^2
  + \binom{p-1}{2}(p+1)^2
  + (p-1) \cdot p^2
  + \frac{p(p-1)}{2}(p-1)^2
  = (p^2-1)(p^2-p)
  = |G|,
\]
consistent with $\sum_\rho d_\rho^2 = |G|$.

\section{The Random Walk Operator via Representations}

For the uniform measure on the symmetric transvection generating set $S_{\mathrm{tv}}$,
the Fourier transform at a representation $\rho$ is
\[
  \hat\mu(\rho) = \frac{1}{4}\bigl(\rho(T_{12}) + \rho(T_{12}^{-1}) + \rho(T_{21}) + \rho(T_{21}^{-1})\bigr).
\]
Since each generator is unipotent, its image $\rho(T_{ij})$ under any
representation is a unipotent matrix.
The operator norm $\|\hat\mu(\rho)\|_{\mathrm{op}}$ governs how fast the
corresponding Fourier mode decays.

For the \textbf{trivial representation} $\mathbf{1}$, $\hat\mu(\mathbf{1}) = 1$.
For all non-trivial representations, the key quantity is the \textbf{spectral
radius} $\beta_\rho = \|\hat\mu(\rho)\|_{\mathrm{op}} < 1$, and the upper bound
lemma gives
\[
  4\|\mu^{*t} - U_G\|_{\mathrm{TV}}^2
  \leq \sum_{\rho \neq \mathbf{1}} d_\rho \cdot d_\rho \cdot \beta_\rho^{2t}
  = \sum_{\rho \neq \mathbf{1}} d_\rho^2\,\beta_\rho^{2t}.
\]
The dominant term as $t \to \infty$ comes from the representation with the
largest $\beta_\rho$, which determines the asymptotic mixing rate.

\section{Example Computations}

\subsection{Exact Analysis for $\GL_2(\mathbb{F}_2) \cong S_3$}

$\GL_2(\mathbb{F}_2)$ has 6 elements and is isomorphic to $S_3$.
Its three irreducible representations are: the trivial rep ($d = 1$), the sign
rep ($d = 1$), and the standard rep ($d = 2$).

For the lazy transvection walk with $\tilde\mu$ uniform on $\{e, T_{12}, T_{21}\}$:
\begin{itemize}
  \item Trivial representation: $\hat{\tilde\mu}(\mathbf{1}) = 1$.
  \item Sign representation: Since $T_{12}$ and $T_{21}$ correspond to
        transpositions in $S_3$ (which are odd permutations),
        $\hat{\tilde\mu}(\mathrm{sgn}) = \frac{1}{3}(1 + (-1) + (-1)) = -\frac{1}{3}$.
  \item Standard representation (dim 2): $\hat{\tilde\mu}(\mathrm{std})$
        has operator norm $\frac{1}{3}$, computed from the explicit matrices.
\end{itemize}

The upper bound lemma gives
\[
  4\|\tilde\mu^{*t} - U_G\|_{\mathrm{TV}}^2
  \leq 1 \cdot \left(\frac{1}{3}\right)^{2t} + 2 \cdot \left(\frac{1}{3}\right)^{2t}
  = 3 \cdot \left(\frac{1}{9}\right)^t,
\]
so $\|\tilde\mu^{*t} - U_G\|_{\mathrm{TV}} \leq \frac{\sqrt{3}}{2} \cdot (1/3)^t$.
This matches the spectral gap computation: $\gamma = 2/3$, $(1-\gamma)^t = (1/3)^t$.

\subsection{Convergence for $\GL_2(\mathbb{F}_3)$}

For $p = 3$, $|G| = 48$ and $|\SL_2(\mathbb{F}_3)| = 24$.
The transvection walk lives on $\SL_2(\mathbb{F}_3)$, while the
transvections+swap walk is ergodic on all of $\GL_2(\mathbb{F}_3)$.

Empirically (from Monte Carlo simulation with $N = 3000$ independent walks):
\begin{itemize}
  \item Transvection walk on $\SL_2(\mathbb{F}_3)$: TVD $< 0.05$ by $t \approx 20$.
  \item Transvections+swap walk on $\GL_2(\mathbb{F}_3)$: TVD $< 0.05$ by $t \approx 25$.
\end{itemize}
The additional swap generator slightly increases the mixing time (larger
generating set means each generator has lower weight), but the walk now covers
the full group.

\subsection{Scaling with $p$}

As $p$ grows, the group order grows as $\sim p^4$ while the generating set size
stays fixed at 4 (or 5 with swap).
The spectral gap $\gamma$ for the transvection generators decreases with $p$,
reflecting the increasing complexity of the group.
Empirical mixing times (TVD $< 0.05$, transvections+swap):

\begin{center}
\begin{tabular}{r r r}
\hline
$p$ & $|\GL_2(\mathbb{F}_p)|$ & approx.\ $t_{\mathrm{mix}}(0.05)$ \\
\hline
$2$ & $6$ & $8$ \\
$3$ & $48$ & $25$ \\
$5$ & $480$ & $\sim 70$ \\
$7$ & $2016$ & $\sim 200$ \\
\hline
\end{tabular}
\end{center}

The mixing time grows faster than $\log|G|$ (which would be optimal) but
appears to grow sub-linearly relative to $|G|$ itself.

\section{Sharp Asymptotics: Kassabov's Theorem}

The precise mixing time for transvection generators on $\GL_n(\mathbb{F}_q)$
(for general $n$ and prime power $q$) was established by Kassabov (2007):

\begin{theorem}[Kassabov 2007]
The random walk on $\GL_n(\mathbb{F}_q)$ driven by elementary matrices
(transvections) has mixing time $\Theta(n^2 \log q)$ in total variation.
For $n = 2$ this gives $\Theta(\log p)$ steps to mix on $\SL_2(\mathbb{F}_p)$.
\end{theorem}

This is \emph{optimal} in the sense that $\Omega(\log |G|)$ steps are always
needed (since TVD is at most 1 and decreases by at most a constant factor per step).
The proof uses the full strength of the representation theory of $\GL_n(\mathbb{F}_q)$
and careful bounds on the Fourier transforms $\hat\mu(\rho)$ for each family of
irreducibles.

\section{Connection to Cayley Expander Graphs}

The random walk on $G$ with generating set $S$ is intimately connected to the
\textbf{Cayley graph} $\mathrm{Cay}(G, S)$.
The spectral gap $\gamma$ of the walk equals the spectral gap of this graph,
and a sequence of Cayley graphs with bounded degree and $\gamma \geq \epsilon > 0$
(independent of $|G|$) forms a family of \textbf{expander graphs}.

Kassabov's theorem shows that $\mathrm{Cay}(\SL_2(\mathbb{F}_p), S_{\mathrm{tv}})$
is an expander family in $p$: the spectral gap is bounded below by a positive
constant independent of $p$.
This has striking consequences in combinatorics and computer science, including
the explicit construction of error-correcting codes and pseudorandom generators.

\section*{References}

\begin{itemize}
  \item P.\ Diaconis and M.\ Shahshahani, \emph{Generating a random permutation with random transpositions}, Z.\ Wahrsch.\ Verw.\ Gebiete 57 (1981), 159--179.
  \item P.\ Diaconis, \emph{Group Representations in Probability and Statistics}, IMS Lecture Notes, 1988.
  \item M.\ Kassabov, \emph{Symmetric groups and expander graphs}, Invent.\ Math.\ 170 (2007), 327--354.
  \item J.-P.\ Serre, \emph{Linear Representations of Finite Groups}, Springer, 1977.
  \item R.\ Steinberg, \emph{The representations of $\GL(3, q)$, $\GL(4, q)$, $PGL(3, q)$, and $PGL(4, q)$}, Canad.\ J.\ Math.\ 3 (1951), 225--235.
\end{itemize}
